\begin{corollary}[Endpoint weak-type bound for a single coordinate $R_{\gamma,j}$]
Let $\mathcal{M}(\gamma_n;\mathbb{R}^n)$ be the space of $\gamma_n$-a.e.\
equivalence classes of measurable $\mathbb{R}^n$-valued functions on
$\mathbb{R}^n$, and let
$T : L^1(\gamma_n;\mathbb{R}) \to \mathcal{M}(\gamma_n;\mathbb{R}^n)$ be any
linear map with the two properties
\begin{enumerate}
\item $T(\iota v) = R_\gamma v$ for every $v \in L^2(\gamma_n)$, with $\iota$ the
inclusion of \noderef{ToL1} and $R_\gamma$ the transform of
\noderef{RieszTransform}; and
\item $t\,\gamma_n\bigl(\{x : |Tf(x)|_{\mathbb{R}^n} > t\}\bigr)
       \leq 2\,\|f\|_{L^1(\gamma_n)}$
for every $f \in L^1(\gamma_n;\mathbb{R})$ and every real $t > 0$.
\end{enumerate}
Then for every $j \in \{1,\ldots,n\}$, every $f \in L^1(\gamma_n;\mathbb{R})$ and
every real $t > 0$, the $j$th coordinate function $x \mapsto (Tf(x))_j$ satisfies
\[
  t\,\gamma_n\bigl(\{x : |(Tf(x))_j| > t\}\bigr) \leq 2\,\|f\|_{L^1(\gamma_n)} .
\]

This is the first assertion of the reference paper's corollary on coordinate and
historical conventions,
$\sup_{t>0} t\,\gamma_n(\{|R_{\gamma,j}f|>t\}) \leq 2\|f\|_1$, stated per level
$t$ rather than as a supremum; the two forms are the same assertion, since a
family of real numbers is bounded by a constant exactly when its supremum is.
It differs from property (2) above, and hence from the endpoint theorem it is
drawn from, in that the distribution function is taken for the \emph{scalar}
function $(Tf)_j$ rather than for the Euclidean norm $|Tf|_{\mathbb{R}^n}$ of the
whole vector; this is a strictly weaker conclusion for each fixed $j$ and is the
form in which the coordinate transforms $R_{\gamma,j}$ are usually quoted.

Properties (1) and (2) are written out here rather than imported by citation.
They are exactly the two properties for which \noderef{EndpointTheorem} asserts
existence and uniqueness of a map $T$, so the universally quantified statement
above is not vacuous: there is such a $T$, there is only one, and the statement
about that unique extension of $R_\gamma$ --- the one the reference paper writes
$R_{\gamma,j}$ for in its $j$th coordinate --- is the case of the present
statement at that $T$.
\end{corollary}
\begin{proof}
Fix $j$, $f$ and $t > 0$, and choose a representative of $Tf$, so that $Tf(x)$ is
a point of $\mathbb{R}^n$ for every $x$ and $(Tf(x))_j$ is its $j$th coordinate.

\emph{Step 1: the pointwise bound.}
For every $x \in \mathbb{R}^n$ the $j$th coordinate of a vector is bounded in
absolute value by its Euclidean norm:
\[
  |(Tf(x))_j| \;=\; \bigl((Tf(x))_j^2\bigr)^{1/2}
    \;\leq\; \Bigl(\sum_{i=1}^n (Tf(x))_i^2\Bigr)^{1/2}
    \;=\; |Tf(x)|_{\mathbb{R}^n},
\]
the inequality because the omitted summands $(Tf(x))_i^2$, $i \neq j$, are
nonnegative and $s \mapsto s^{1/2}$ is nondecreasing on $[0,\infty)$.  This is the
pointwise inequality $|R_{\gamma,j}f| \leq |R_\gamma f|_{\mathbb{R}^n}$ with which
the reference paper begins its proof.

\emph{Step 2: inclusion of the level sets.}
If $t < |(Tf(x))_j|$ then $t < |Tf(x)|_{\mathbb{R}^n}$ by Step 1, so
\[
  \{x : |(Tf(x))_j| > t\} \subseteq \{x : |Tf(x)|_{\mathbb{R}^n} > t\} .
\]
By monotonicity of a measure (shared context),
$\gamma_n(\{|(Tf)_j| > t\}) \leq \gamma_n(\{|Tf|_{\mathbb{R}^n} > t\})$.

\emph{Step 3: conclusion.}
Multiplying the inequality of Step 2 by $t > 0$ and then applying hypothesis (2)
at this $f$ and this $t$,
\[
  t\,\gamma_n\bigl(\{|(Tf)_j| > t\}\bigr)
    \;\leq\; t\,\gamma_n\bigl(\{|Tf|_{\mathbb{R}^n} > t\}\bigr)
    \;\leq\; 2\,\|f\|_{L^1(\gamma_n)} ,
\]
which is the assertion.  Both quantities are elements of $[0,+\infty]$ and the
first inequality is multiplication of an inequality in $[0,+\infty]$ by the
nonnegative constant $t$, which preserves the order.
\end{proof}
